\chapter{文件系统}
\label{CH:FS}
%

文件系统的目的是组织和存储数据。文件系统通常支持用户和应用程序之间的数据共享，以及
\indextext{持久性（persistence）}，
以便数据在重启后仍然可用。

xv6 文件系统提供类似 Unix 的文件、目录和路径名
（参见第~\ref{CH:UNIX} 章），并将其数据存储在 virtio 磁盘上以实现持久性。文件系统解决了几个挑战：
\begin{itemize}

\item 文件系统需要磁盘上的数据结构来表示命名的目录和文件树，
记录保存每个文件内容的块的标识，并记录磁盘的哪些区域是空闲的。
\item 文件系统必须支持
\indextext{崩溃恢复（crash recovery）}。
也就是说，如果发生崩溃（例如，电源故障），文件系统必须在重启后仍然正常工作。
风险在于崩溃可能会中断一系列更新，导致磁盘上的数据结构不一致（例如，一个块既在文件中使用又被标记为空闲）。
\item 不同的进程可能同时对文件系统进行操作，
因此文件系统代码必须协调以保持不变量。
\item 访问磁盘比访问内存慢几个数量级，
因此文件系统必须维护一个热门块的内存缓存。

\end{itemize}

本章的其余部分解释了 xv6 如何解决这些挑战。
%%
%%  -------------------------------------------
%%
\section{概述}

xv6 文件系统实现分为七个层次，如图~\ref{fig:fslayer} 所示。
磁盘层在 virtio 硬盘上读写块。
缓冲区缓存层缓存磁盘块并同步对它们的访问，
确保一次只有一个内核进程可以修改存储在任何特定块中的数据。
日志层允许更高层将对多个块的更新包装在一个
\indextext{事务（transaction）}中，
并确保在崩溃面前原子地更新块（即，全部更新或都不更新）。
索引节点层提供单个文件，每个文件由一个具有唯一 i-number 的
\indextext{索引节点（inode）}
和一些保存文件数据的块表示。
目录层将每个目录实现为一种特殊的 inode，其内容是一系列目录项，每个目录项包含文件的名称和 i-number。
路径名层提供分层路径名，如
\lstinline{/usr/rtm/xv6/fs.c}，
并通过递归查找来解析它们。
文件描述符层使用文件系统接口抽象许多 Unix 资源（例如，管道、设备、文件等），简化应用程序程序员的工作。

\begin{figure}[t]
\centering
\includegraphics[scale=0.5]{fig/fslayer.pdf}
\caption{xv6 文件系统的层次。}
\label{fig:fslayer}
\end{figure}

磁盘硬件传统上将磁盘上的数据表示为编号序列的 512 字节
\textit{块}
\index{块（block）}
（也称为
\textit{扇区}）：
\index{扇区（sector）}
扇区 0 是前 512 字节，扇区 1 是接下来的 512 字节，依此类推。操作系统用于其文件系统的块大小可能与磁盘使用的扇区大小不同，但通常块大小是扇区大小的倍数。Xv6 将已读入内存的块副本保存在类型为
\lstinline{struct buf}
\lineref{kernel/buf.h:/^struct.buf/}
的对象中。
该结构中存储的数据有时与磁盘不同步：它可能尚未从磁盘读入（磁盘正在处理但尚未返回扇区的内容），或者它可能已被软件更新但尚未写入磁盘。

文件系统必须有一个计划来确定它在磁盘上存储 inode 和内容块的位置。
为此，xv6 将磁盘分成几个部分，如图~\ref{fig:fslayout} 所示。
文件系统不使用块 0（它保存引导扇区）。
块 1 称为
\indextext{超级块（superblock）}；
它包含文件系统的元数据（文件系统大小（以块为单位）、数据块数量、inode 数量和日志中的块数）。
从块 2 开始的块保存日志。日志之后是 inode，每个块包含多个 inode。
之后是位图块，跟踪哪些数据块正在使用。
剩余的块是数据块；每个块要么在位图块中标记为空闲，要么保存文件或目录的内容。
超级块由一个称为
\indexcode{mkfs}
的单独程序填充，该程序构建初始文件系统。

本章的其余部分讨论每一层，从缓冲区缓存开始。
注意低层中精心选择的抽象如何简化高层设计的情况。
%%
%%  -------------------------------------------
%%
\section{缓冲区缓存层}
\label{s:bcache}

缓冲区缓存有两个工作：(1) 同步对磁盘块的访问，以确保内存中只有一个块的副本，并且一次只有一个内核线程使用该副本；
(2) 缓存热门块，以便不需要从慢速磁盘重新读取。
代码在
\lstinline{bio.c} 中。

缓冲区缓存导出的主要接口包括
\indexcode{bread}
和
\indexcode{bwrite}；
前者获取一个包含块副本的
\indextext{缓冲区（buf）}，可以在内存中读取或修改，后者将修改后的缓冲区写入磁盘上的适当块。
内核线程在使用完缓冲区后必须通过调用
\indexcode{brelse}
释放它。
缓冲区缓存使用每个缓冲区的睡眠锁来确保一次只有一个线程使用每个缓冲区
（因此每个磁盘块）；
\lstinline{bread}
返回一个锁定的缓冲区，
\lstinline{brelse}
释放锁。

让我们回到缓冲区缓存。
缓冲区缓存有固定数量的缓冲区来保存磁盘块，
这意味着如果文件系统请求一个尚未在缓存中的块，缓冲区缓存必须回收当前保存其他块的缓冲区。
缓冲区缓存为新块回收最近最少使用的缓冲区。
假设是最近最少使用的缓冲区是最不可能很快再次使用的缓冲区。

\begin{figure}[t]
\centering
\includegraphics[scale=0.5]{fig/fslayout.pdf}
\caption{xv6 文件系统的结构。}
\label{fig:fslayout}
\end{figure}
%%
%%  -------------------------------------------
%%
\section{代码：缓冲区缓存}

缓冲区缓存是缓冲区的双向链表。
函数
\indexcode{binit}，
由
\indexcode{main}
\lineref{kernel/main.c:/binit/}
调用，
用静态数组
\lstinline{buf}
中的
\indexcode{NBUF}
个缓冲区初始化列表
\linerefs{kernel/bio.c:/Create.linked.list/,/^..}/}。
所有其他对缓冲区缓存的访问都通过
\indexcode{bcache.head}
引用链表，而不是
\lstinline{buf}
数组。

缓冲区有两个与之关联的状态字段。
字段
\indexcode{valid}
表示缓冲区包含块的副本。
字段 \indexcode{disk}
表示缓冲区内容已交给磁盘，磁盘可能会更改缓冲区（例如，将数据从磁盘写入 \lstinline{data}）。

\lstinline{bread}
\lineref{kernel/bio.c:/^bread/}
调用
\indexcode{bget}
获取给定扇区的缓冲区
\lineref{kernel/bio.c:/b.=.bget/}。
如果需要从磁盘读取缓冲区，
\lstinline{bread}
调用
\indexcode{virtio_disk_rw}
在返回缓冲区之前执行此操作。

\lstinline{bget}
\lineref{kernel/bio.c:/^bget/}
扫描缓冲区列表以查找具有给定设备和扇区号的缓冲区
\linerefs{kernel/bio.c:/Is.the.block.already/,/^..}/}。
如果有这样的缓冲区，
\indexcode{bget}
获取缓冲区的睡眠锁。
\lstinline{bget}
然后返回锁定的缓冲区。

如果没有给定扇区的缓存缓冲区，
\indexcode{bget}
必须创建一个，可能重用保存不同扇区的缓冲区。
它第二次扫描缓冲区列表，查找未使用的缓冲区（\lstinline{b->refcnt = 0}）；
任何这样的缓冲区都可以使用。
\lstinline{bget}
编辑缓冲区元数据以记录新的设备和扇区号，并获取其睡眠锁。
注意赋值
\lstinline{b->valid = 0}
确保
\lstinline{bread}
将从磁盘读取块数据，
而不是错误地使用缓冲区的先前内容。

每个磁盘扇区最多有一个缓存缓冲区很重要，以确保读取者看到写入，
而且因为文件系统使用缓冲区上的锁进行同步。
\lstinline{bget}
通过从第一个循环检查块是否已缓存到第二个循环声明块现在已缓存（通过设置
\lstinline{dev}、
\lstinline{blockno}
和
\lstinline{refcnt}）持续持有
\lstinline{bcache.lock}
来确保这个不变量。
这使得检查块是否存在以及（如果不存在）指定一个缓冲区来保存块是原子的。

\lstinline{bget}
在
\lstinline{bcache.lock}
临界区外获取缓冲区的睡眠锁是安全的，
因为非零的
\lstinline{b->refcnt}
阻止缓冲区被重用于不同的磁盘块。
睡眠锁保护块缓冲内容的读写，而
\lstinline{bcache.lock}
保护关于哪些块被缓存的信息。

如果所有缓冲区都忙，那么同时执行文件系统调用的进程太多了；
\lstinline{bget}
发生 panic。
更优雅的响应可能是睡眠直到缓冲区空闲，
尽管那样会有死锁的可能性。

一旦
\indexcode{bread}
读取了磁盘（如果需要）并将缓冲区返回给其调用者，调用者就拥有缓冲区的独占使用权，可以读取或写入数据字节。
如果调用者确实修改了缓冲区，它必须在释放缓冲区之前调用
\indexcode{bwrite}
将更改的数据写入磁盘。
\lstinline{bwrite}
\lineref{kernel/bio.c:/^bwrite/}
调用
\indexcode{virtio_disk_rw}
与磁盘硬件通信。

当调用者使用完缓冲区时，它必须调用
\indexcode{brelse}
释放它。
（名称
\lstinline{brelse}，
b-release 的缩写，
是神秘的但值得学习的：
它起源于 Unix，也用于 BSD、Linux 和 Solaris。）
\lstinline{brelse}
\lineref{kernel/bio.c:/^brelse/}
释放睡眠锁并将缓冲区移到链表的前端
\linerefs{kernel/bio.c:/b->next->prev.=.b->prev/,/bcache.head.next.=.b/}。
移动缓冲区使列表按缓冲区最近使用的时间（意味着释放）排序：
列表中的第一个缓冲区是最近使用的，
最后一个是最近最少使用的。
\lstinline{bget}
中的两个循环利用这一点：
扫描现有缓冲区的最坏情况下必须处理整个列表，
但先检查最近使用的缓冲区（从
\lstinline{bcache.head}
开始并跟随
\lstinline{next}
指针）将在引用局部性良好时减少扫描时间。
选择要重用的缓冲区的扫描通过后向扫描（跟随
\lstinline{prev}
指针）选择最近最少使用的缓冲区。
%%
%%  -------------------------------------------
%%
\section{日志层}

文件系统设计中最有趣的问题之一是崩溃恢复。
问题出现是因为许多文件系统操作涉及对磁盘的多次写入，而在部分写入后发生崩溃可能会使磁盘上的文件系统处于不一致状态。
例如，假设在文件截断期间发生崩溃（将文件长度设置为零并释放其内容块）。
根据磁盘写入的顺序，崩溃可能留下一个引用标记为空闲的内容块的 inode，
或者可能留下一个已分配但未引用的内容块。

后者相对无害，但引用已释放块的 inode 可能在重启后导致严重问题。
重启后，内核可能将该块分配给另一个文件，现在我们有两个不同的文件无意中指向同一个块。
如果 xv6 支持多用户，这种情况可能是安全问题，因为旧文件的所有者将能够读写新文件中的块，而新文件由不同的用户拥有。

Xv6 使用一种简单的日志形式来解决文件系统操作期间崩溃的问题。
xv6 系统调用不直接写入磁盘上的文件系统数据结构。
相反，它将它希望进行的所有磁盘写入的描述放在磁盘上的
\indextext{日志（log）}
中。
一旦系统调用记录了其所有写入，它就向磁盘写入一个特殊的
\indextext{提交（commit）}
记录，指示日志包含一个完整的操作。
此时，系统调用将写入复制到磁盘上的文件系统数据结构。
在这些写入完成后，系统调用擦除磁盘上的日志。

如果系统崩溃并重启，文件系统代码在运行任何进程之前按如下方式从崩溃中恢复。
如果日志被标记为包含完整操作，则恢复代码将写入复制到磁盘上文件系统中的适当位置。
如果日志未被标记为包含完整操作，恢复代码将忽略日志。
恢复代码通过擦除日志完成恢复。

为什么 xv6 的日志能解决文件系统操作期间崩溃的问题？
如果崩溃发生在操作提交之前，则磁盘上的日志不会被标记为完整，恢复代码将忽略它，磁盘的状态就像操作甚至没有开始一样。
如果崩溃发生在操作提交之后，则恢复将重放操作的所有写入，如果操作已经开始将它们写入磁盘数据结构，则可能重复它们。
无论哪种情况，日志都使操作相对于崩溃是原子的：恢复后，操作的所有写入要么都出现在磁盘上，要么都不出现。
%%
%%
%%
\section{日志设计}

日志位于已知的固定位置，在超级块中指定。
它由一个头块和一系列更新的块副本（``日志块''）组成。
头块包含一个扇区号数组，每个日志块一个，以及日志块的数量。
磁盘上头块中的计数要么为零，表示日志中没有事务，
要么为非零，表示日志包含具有指示数量的日志块的完整提交事务。
Xv6 在事务提交时写入头块，但之前不写入，并在将日志块复制到文件系统后将计数设置为零。
因此，事务中途崩溃将导致日志头块中的计数为零；
提交后崩溃将导致计数为非零。

每个系统调用的代码指示必须相对于崩溃原子的写入序列的开始和结束。
为了允许不同进程并发执行文件系统操作，
日志系统可以将多个系统调用的写入累积到一个事务中。
因此，单个提交可能涉及多个完整系统调用的写入。
为了避免将系统调用跨事务分割，日志系统仅在没有文件系统系统调用正在进行时才提交。

将多个事务一起提交的想法称为
\indextext{组提交（group commit）}。
组提交减少了磁盘操作的数量，因为它将提交的固定成本分摊到多个操作上。
组提交还同时向磁盘系统提交更多并发写入，
可能允许磁盘在单次磁盘旋转期间写入它们所有。
Xv6 的 virtio 驱动程序不支持这种
\indextext{批处理（batching）}，
但 xv6 的文件系统设计允许这样做。

Xv6 在磁盘上分配固定量的空间来保存日志。
事务中系统调用写入的块总数必须适合该空间。
这有两个后果。
不能允许单个系统调用写入比日志空间更多的不同块。
这对大多数系统调用来说不是问题，但有两个可能写入许多块：
\indexcode{write}
和
\indexcode{unlink}。
大文件写入可能写入许多数据块和许多位图块以及一个 inode 块；
取消链接大文件可能写入许多位图块和一个 inode。
Xv6 的写入系统调用将大写入拆分为多个适合日志的较小写入，
而
\lstinline{unlink}
不会引起问题，因为实际上 xv6 文件系统只使用一个位图块。
有限日志空间的另一个后果是日志系统不能允许系统调用开始，除非它确定系统调用的写入将适合日志中剩余的空间。
%%
%%
%%
\section{代码：日志}

系统调用中日志的典型使用如下所示：
\begin{lstlisting}[]
  begin_op();
  ...
  bp = bread(...);
  bp->data[...] = ...;
  log_write(bp);
  ...
  end_op();
\end{lstlisting}

\indexcode{begin_op}
\lineref{kernel/log.c:/^begin.op/}
等待直到日志系统当前未提交，并且有足够的未预留日志空间来容纳此调用的写入。
\lstinline{log.outstanding}
计数已预留日志空间的系统调用数量；
总预留空间是
\lstinline{log.outstanding}
乘以
\lstinline{MAXOPBLOCKS}。
递增
\lstinline{log.outstanding}
既预留空间又防止在此系统调用期间发生提交。
代码保守地假设每个系统调用可能写入多达
\lstinline{MAXOPBLOCKS}
个不同的块。

\indexcode{log_write}
\lineref{kernel/log.c:/^log.write/}
充当
\indexcode{bwrite}
的代理。
它在内存中记录块的扇区号，
在磁盘上的日志中为它预留一个槽位，
并将缓冲区固定在块缓存中以防止块缓存逐出它。
该块必须保持在缓存中直到提交：
在此之前，缓存的副本是修改的唯一记录；
它不能在提交后写入其在磁盘上的位置；
而且同一事务中的其他读取必须看到修改。
\lstinline{log_write}
注意到当一个块在单个事务中被多次写入时，并为该块在日志中分配相同的槽位。
这种优化通常称为
\indextext{吸收（absorption）}。
常见的情况是，例如，包含多个文件的 inode 的磁盘块在事务内被多次写入。
通过将多个磁盘写入吸收到一个中，文件系统可以节省日志空间，并且可以实现更好的性能，因为只有一个磁盘块副本必须写入磁盘。

\indexcode{end_op}
\lineref{kernel/log.c:/^end.op/}
首先递减未完成系统调用的计数。
如果计数现在为零，它通过调用
\lstinline{commit()}
提交当前事务。
此过程有四个阶段。
\lstinline{write_log()}
\lineref{kernel/log.c:/^write.log/}
将事务中修改的每个块从缓冲区缓存复制到磁盘上日志中的其槽位。
\lstinline{write_head()}
\lineref{kernel/log.c:/^write.head/}
将头块写入磁盘：这是提交点，写入后崩溃将导致恢复从事务的日志重放写入。
\indexcode{install_trans}
\lineref{kernel/log.c:/^install_trans/}
从日志中读取每个块并将其写入文件系统中的适当位置。
最后
\lstinline{end_op}
写入计数为零的日志头；
这必须在下一个事务开始写入日志块之前发生，以便崩溃不会导致恢复使用一个事务的头与后续事务的日志块。

\indexcode{recover_from_log}
\lineref{kernel/log.c:/^recover_from_log/}
由
\indexcode{initlog}
\lineref{kernel/log.c:/^initlog/}
调用，
后者在第一个用户进程运行之前从 \indexcode{fsinit}\lineref{kernel/fs.c:/^fsinit/} 在引导期间调用
\lineref{kernel/proc.c:/fsinit/}。
它读取日志头，如果头指示日志包含已提交的事务，则模拟
\lstinline{end_op}
的操作。

日志的一个示例使用发生在
\indexcode{filewrite}
\lineref{kernel/file.c:/^filewrite/}。
事务如下所示：
\begin{lstlisting}[]
      begin_op();
      ilock(f->ip);
      r = writei(f->ip, ...);
      iunlock(f->ip);
      end_op();
\end{lstlisting}
这段代码包装在一个循环中，将大写入拆分为每次只有几个扇区的单独事务，以避免溢出日志。
对
\indexcode{writei}
的调用作为此事务的一部分写入许多块：文件的 inode、一个或多个位图块和一些数据块。
%%
%%
%%
\section{代码：块分配器}

文件和目录内容存储在磁盘块中，必须从空闲池中分配。
Xv6 的块分配器在磁盘上维护一个空闲位图，每个块一位。
零位表示相应的块是空闲的；
一位表示它正在使用。
程序
\lstinline{mkfs}
设置对应于引导扇区、超级块、日志块、inode 块和位图块的位。

块分配器提供两个函数：
\indexcode{balloc}
分配一个新的磁盘块，
\indexcode{bfree}
释放一个块。
\lstinline{balloc}
中的循环在
\lineref{kernel/fs.c:/^..for.b.=.0/}
考虑从块 0 到
\lstinline{sb.size}
（文件系统中的块数）的每个块。
它寻找位图为零的块，表示它是空闲的。
如果
\lstinline{balloc}
找到这样的块，它更新位图并返回该块。
为了提高效率，循环分为两部分。
外循环读取每个位图块。
内循环检查单个位图块中的所有
Bits-Per-Block (\lstinline{BPB})
位。
如果两个进程试图同时分配一个块可能发生的竞争被缓冲区缓存只允许一个进程一次使用任何一个位图块的事实阻止。

\lstinline{bfree}
\lineref{kernel/fs.c:/^bfree/}
找到正确的位图块并清除正确的位。
同样，
\lstinline{bread}
和
\lstinline{brelse}
暗示的独占使用避免了对显式锁定的需要。

与本章其余部分描述的大部分代码一样，
\lstinline{balloc}
和
\lstinline{bfree}
必须在事务内调用。
%%
%%  -------------------------------------------
%%
\section{索引节点层}

术语
\indextext{索引节点（inode）}
可以有两种相关含义之一。
它可能指包含文件大小和数据块号列表的磁盘上的数据结构。
或者``inode''可能指内存中的 inode，它包含磁盘上 inode 的副本以及内核内需要的额外信息。

磁盘上的 inode 被打包到磁盘上称为 inode 块的连续区域中。
每个 inode 大小相同，因此给定一个数字 n，很容易在磁盘上找到第 n 个 inode。
事实上，这个数字 n，称为 inode 号或 i-number，是 inode 在实现中的标识方式。

磁盘上的 inode 由
\indexcode{struct dinode}
\lineref{kernel/fs.h:/^struct.dinode/}
定义。
\lstinline{type}
字段区分文件、目录和特殊文件（设备）。
类型为零表示磁盘上的 inode 是空闲的。
\lstinline{nlink}
字段计算引用此 inode 的目录条目数，以识别何时应释放磁盘上的 inode 及其数据块。
\lstinline{size}
字段记录文件中内容的字节数。
\lstinline{addrs}
数组记录保存文件内容的磁盘块的块号。

内核在称为 \indexcode{itable} 的表中将活动 inode 集保存在内存中；
\indexcode{struct inode}
\lineref{kernel/file.h:/^struct.inode/}
是磁盘上
\lstinline{struct}
\lstinline{dinode}
的内存副本。
内核仅在内存中有 C 指针引用该 inode 时才将 inode 存储在内存中。
\lstinline{ref}
字段计算引用内存中 inode 的 C 指针数量，如果引用计数降为零，内核将从内存中丢弃该 inode。
\indexcode{iget}
和
\indexcode{iput}
函数获取和释放指向 inode 的指针，修改引用计数。
指向 inode 的指针可能来自文件描述符、当前工作目录和瞬态内核代码，如
\lstinline{kexec}。

xv6 的 inode 代码中有四种锁或类似锁的机制。
\lstinline{itable.lock}
保护 inode 在 inode 表中最多出现一次的不变量，以及内存中 inode 的
\lstinline{ref}
字段计算内存中指向该 inode 的指针数量的不变量。
每个内存中的 inode 有一个包含睡眠锁的
\lstinline{lock}
字段，它确保对 inode 字段（如文件长度）以及 inode 的文件或目录内容块的独占访问。
如果 inode 的
\lstinline{ref}
大于零，它会导致系统维护表中的 inode，而不是为不同的 inode 重用表条目。
最后，每个 inode 包含一个
\lstinline{nlink}
字段（在磁盘上，如果在内存中则复制在内存中），计算引用文件的目录条目数；
如果其链接计数大于零，xv6 不会释放 inode。

\lstinline{iget()}
返回的
\lstinline{struct}
\lstinline{inode}
指针保证在相应调用
\lstinline{iput()}
之前有效；
inode 不会被删除，指针引用的内存不会为不同的 inode 重用。
\lstinline{iget()}
提供对 inode 的非独占访问，因此可以有多个指针指向同一个 inode。
文件系统代码的许多部分依赖于
\lstinline{iget()}
的这种行为，
既为了持有对 inode 的长期引用（作为打开的文件和当前目录），也为了在操作多个 inode 的代码中防止竞争同时避免死锁（如路径名查找）。

\lstinline{iget}
返回的
\lstinline{struct}
\lstinline{inode}
可能没有任何有用的内容。
为了确保它持有磁盘上 inode 的副本，代码必须调用
\indexcode{ilock}。
这会锁定 inode（以便没有其他进程可以
\lstinline{ilock}
它）并从磁盘读取 inode，如果尚未读取的话。
\lstinline{iunlock}
释放 inode 上的锁。
将 inode 指针的获取与锁定分开有助于在某些情况下避免死锁，例如在目录查找期间。
多个进程可以持有
\lstinline{iget}
返回的 inode 的 C 指针，
但一次只有一个进程可以锁定 inode。

inode 表仅存储内核代码或数据结构持有 C 指针的 inode。
它的主要工作是同步多个进程的访问。
inode 表也恰好缓存经常使用的 inode，但缓存是次要的；
如果 inode 被频繁使用，缓冲区缓存可能会将其保留在内存中。
修改内存中 inode 的代码使用
\lstinline{iupdate}
将其写入磁盘。
%%
%%  -------------------------------------------
%%
\section{代码：索引节点}

要分配一个新的 inode（例如，创建文件时），
xv6 调用
\indexcode{ialloc}
\lineref{kernel/fs.c:/^ialloc/}。
\lstinline{ialloc}
类似于
\indexcode{balloc}：
它遍历磁盘上的 inode 结构，一次一个块，寻找一个标记为空闲的 inode。
当它找到一个时，它通过将新的
\lstinline{type}
写入磁盘来声明它，然后通过尾调用返回 inode 表中的一个条目
\indexcode{iget}
\lineref{kernel/fs.c:/return.iget\(dev..inum\)/}。
\lstinline{ialloc}
的正确操作依赖于这样一个事实：一次只有一个进程可以持有对
\lstinline{bp}
的引用：
\lstinline{ialloc}
可以确信没有其他进程同时看到该 inode 可用并试图声明它。

\lstinline{iget}
\lineref{kernel/fs.c:/^iget/}
查看 inode 表以查找具有所需设备和 inode 号的活动条目（\lstinline{ip->ref}
\lstinline{>}
\lstinline{0}）。
如果找到一个，它返回对该 inode 的新引用
\linerefs{kernel/fs.c:/^....if.ip->ref.>.0/,/^....}/}。
当
\indexcode{iget}
扫描时，它记录第一个空槽的位置
\linerefs{kernel/fs.c:/^....if.empty.==.0/,/empty.=.ip/}，
如果需要分配表条目，则使用它。

代码必须在读取或写入其元数据或内容之前使用
\indexcode{ilock}
锁定 inode。
\lstinline{ilock}
\lineref{kernel/fs.c:/^ilock/}
为此目的使用睡眠锁。
一旦
\indexcode{ilock}
获得对 inode 的独占访问，如果需要，它就从磁盘（更可能是缓冲区缓存）读取 inode。
函数
\indexcode{iunlock}
\lineref{kernel/fs.c:/^iunlock/}
释放睡眠锁，
这可能导致任何睡眠的进程被唤醒。

\lstinline{iput}
\lineref{kernel/fs.c:/^iput/}
通过递减引用计数来释放指向 inode 的 C 指针
\lineref{kernel/fs.c:/^..ip->ref--/}。
如果这是最后一个引用，inode 在 inode 表中的槽位现在空闲，可以重用于不同的 inode。

如果
\indexcode{iput}
看到没有 C 指针引用某个 inode 且该 inode 没有链接到它（不出现在任何目录中），则必须释放该 inode 及其数据块。
\lstinline{iput}
调用
\indexcode{itrunc}
将文件截断为零字节，释放数据块；
将 inode 类型设置为 0（未分配）；
并将 inode 写入磁盘
\lineref{kernel/fs.c:/inode.has.no.links.and/}。

\indexcode{iput}
在释放 inode 的情况下的锁定协议值得仔细研究。
一个危险是并发线程可能在
\lstinline{ilock}
中等待使用此 inode（例如，读取文件或列出目录），
并且不会准备发现该 inode 不再被分配。
这不可能发生，因为如果系统调用无法获得指向内存中 inode 的指针，如果它没有链接到它且
\lstinline{ip->ref}
为一，则无法获得指向内存中 inode 的指针。
该引用是调用
\lstinline{iput}
的线程拥有的引用。
另一个主要危险是对
\lstinline{ialloc}
的并发调用可能选择
\lstinline{iput}
正在释放的相同 inode。
这可能发生在
\lstinline{iupdate}
写入磁盘使 inode 具有零类型之后。
这种竞争是良性的；分配线程将礼貌地等待获取 inode 的睡眠锁，然后再读取或写入 inode，此时
\lstinline{iput}
已完成。

\lstinline{iput()}
可以写入磁盘。这意味着使用文件系统的任何系统调用都可能写入磁盘，因为系统调用可能是最后一个引用文件的系统调用。
即使是像
\lstinline{read()}
这样看似只读的调用，也可能最终调用
\lstinline{iput()}。
反过来，这意味着即使是只读的系统调用，如果使用文件系统，也必须包装在事务中。

\lstinline{iput()}
和崩溃之间存在一个具有挑战性的交互。
\lstinline{iput()}
不会立即截断文件，当文件的链接计数降为零时，因为某些进程可能仍持有对内存中 inode 的引用：进程可能仍在读写文件，因为它成功打开了它。
但是，如果在最后一个进程关闭文件的文件描述符之前发生崩溃，则文件将在磁盘上标记为已分配，但没有目录条目指向它。

文件系统以两种方式之一处理这种情况。简单的解决方案是在恢复期间，重启后，文件系统扫描整个文件系统以查找标记为已分配但没有目录条目指向它们的文件。
如果存在任何此类文件，则可以释放它们。

第二种解决方案不需要扫描文件系统。在这种解决方案中，文件系统记录链接计数降为零但引用计数不为零的文件的 inode inumber（例如，在超级块中）。
如果文件系统在其引用计数达到 0 时删除该文件，则通过从列表中删除该 inode 来更新磁盘上的列表。在恢复时，文件系统释放列表中的任何文件。

Xv6 没有实现任一解决方案，这意味着 inode 可能在磁盘上标记为已分配，即使它们不再使用。
这意味着随着时间的推移，xv6 有磁盘空间耗尽的风险。
%%
%%
%%
\section{代码：索引节点内容}

\begin{figure}[t]
\centering
\includegraphics[scale=0.5]{fig/inode.pdf}
\caption{磁盘上文件的表示。}
\label{fig:inode}
\end{figure}

磁盘上的 inode 结构，
\indexcode{struct dinode}，
包含一个大小和一个块号数组（参见图~\ref{fig:inode}）。
inode 数据位于
\lstinline{dinode}
的
\lstinline{addrs}
数组中列出的块中。
前
\indexcode{NDIRECT}
个数据块列在数组的前
\lstinline{NDIRECT}
个条目中；这些块称为
\indextext{直接块（direct blocks）}。
接下来的
\indexcode{NINDIRECT}
个数据块不在 inode 中列出，而是在一个称为
\indextext{间接块（indirect block）}
的数据块中列出。
\lstinline{addrs}
数组中的最后一个条目给出间接块的地址。
因此，文件的前 12 kB（
\lstinline{NDIRECT}
\lstinline{x}
\indexcode{BSIZE}）
字节可以从 inode 中列出的块加载，
而接下来的
\lstinline{256} kB（
\lstinline{NINDIRECT}
\lstinline{x}
\lstinline{BSIZE}）
字节只能在查阅间接块后加载。
这是一种很好的磁盘上表示，但对客户端来说很复杂。
函数
\indexcode{bmap}
管理这种表示，以便更高级别的例程，如
\indexcode{readi}
和
\indexcode{writei}，
我们稍后会看到，不需要管理这种复杂性。
\lstinline{bmap}
返回 inode
\lstinline{ip}
的第
\lstinline{bn}
个数据块的磁盘块号。
如果
\lstinline{ip}
还没有这样的块，
\lstinline{bmap}
分配一个。

函数
\indexcode{bmap}
\lineref{kernel/fs.c:/^bmap/}
从简单的情况开始：前
\indexcode{NDIRECT}
个块在 inode 本身中列出
\linerefs{kernel/fs.c:/^..if.bn.<.NDIRECT/,/^..}/}。
接下来的
\indexcode{NINDIRECT}
个块列在
\lstinline{ip->addrs[NDIRECT]}
的间接块中。
\lstinline{bmap}
读取间接块
\lineref{kernel/fs.c:/bp.=.bread.ip->dev..addr/}
然后从块中的正确位置读取块号
\lineref{kernel/fs.c:/a.=..uint\*.bp->data/}。
如果块号超过
\lstinline{NDIRECT+NINDIRECT}，
\lstinline{bmap}
发生 panic；
\lstinline{writei}
包含防止这种情况的检查
\lineref{kernel/fs.c:/off...n...MAXFILE.BSIZE/}。

\lstinline{bmap}
根据需要分配块。
\lstinline{ip->addrs[]}
或间接条目为零表示未分配块。
当
\lstinline{bmap}
遇到零时，它用新分配的块号替换它们
\linerefs{kernel/fs.c:/^....if..addr.=.*==.0/,/./}
\linerefs{kernel/fs.c:/^....if..addr.*NDIRECT.*==.0/,/./}。

\indexcode{itrunc}
释放文件的块，将 inode 的大小重置为零。
\lstinline{itrunc}
\lineref{kernel/fs.c:/^itrunc/}
从释放直接块开始
\linerefs{kernel/fs.c:/^..for.i.=.0.*NDIRECT/,/^..}/}，
然后是间接块中列出的那些
\linerefs{kernel/fs.c:/^....for.j.=.0.*NINDIRECT/,/^....}/}，
最后是间接块本身
\linerefs{kernel/fs.c:/^....bfree.*NDIRECT/,/./}。

\lstinline{bmap}
使
\indexcode{readi}
和
\indexcode{writei}
很容易访问 inode 的数据。
\lstinline{readi}
\lineref{kernel/fs.c:/^readi/}
首先确保偏移量和计数不超过文件末尾。
从文件末尾之后开始的读取返回错误
\linerefs{kernel/fs.c:/^..if.off.>.ip->size/,/./}
而从文件末尾开始或跨越文件末尾的读取返回比请求更少的字节
\linerefs{kernel/fs.c:/^..if.off.\+.n.>.ip->size/,/./}。
主循环处理文件的每个块，将数据从缓冲区复制到
\lstinline{dst}
\linerefs{kernel/fs.c:/^..for.tot=0/,/^..}/}。
%%  注意：很难为 writei 写行引用，因为很多行与 readi 中的相同。幸运的是，相同的行可能不需要注释。
\indexcode{writei}
\lineref{kernel/fs.c:/^writei/}
与
\indexcode{readi}
相同，但有三个例外：
从文件末尾开始或跨越文件末尾的写入会增长文件，直到最大文件大小
\linerefs{kernel/fs.c:/^..if.off.\+.n.>.MAXFILE/,/./}；
循环将数据复制到缓冲区中而不是从缓冲区复制出来
\lineref{kernel/fs.c:/either.copyin.*bp->data/}；
而且如果写入扩展了文件，
\indexcode{writei}
必须更新其大小
\linerefs{kernel/fs.c:/^..if.off.>.ip->size\)/,/./}。

函数
\indexcode{stati}
\lineref{kernel/fs.c:/^stati\(/}
将 inode 元数据复制到
\lstinline{stat}
结构体中，该结构体通过
\indexcode{stat}
系统调用暴露给用户程序。
%%
%%
%%
\section{代码：目录层}

目录在内部实现上很像文件。
它的 inode 具有类型
\indexcode{T_DIR}
其数据是一系列目录条目。
每个条目是一个
\indexcode{struct dirent}
\lineref{kernel/fs.h:/^struct.dirent/}，
包含一个名称和一个 inode 号。
名称最多为
\indexcode{DIRSIZ}
（14）个字符；
如果更短，则以 NULL（0）字节终止。
inode 号为零的目录条目是空闲的。

函数
\indexcode{dirlookup}
\lineref{kernel/fs.c:/^dirlookup/}
在目录中搜索具有给定名称的条目。
如果找到一个，它返回指向相应 inode 的指针（未锁定），
并设置
\lstinline{*poff}
为目录中条目的字节偏移量，
以防调用者希望编辑它。
如果
\lstinline{dirlookup}
找到具有正确名称的条目，
它更新
\lstinline{*poff}
并通过
\indexcode{iget}
返回一个未锁定的 inode。
\lstinline{dirlookup}
是
\lstinline{iget}
返回未锁定 inode 的原因。
调用者已锁定
\lstinline{dp}，
因此如果查找是针对
\indexcode{.}，
当前目录的别名，
尝试在返回之前锁定 inode 将试图重新锁定
\lstinline{dp}
并导致死锁。
（涉及多个进程和
\indexcode{..}，
父目录的别名；有更复杂的死锁场景；
\lstinline{.}
不是唯一的问题。）
调用者可以解锁
\lstinline{dp}
然后锁定
\lstinline{ip}，
确保它一次只持有一个锁。

函数
\indexcode{dirlink}
\lineref{kernel/fs.c:/^dirlink/}
将给定的名称和 inode 号的新目录条目写入目录
\lstinline{dp}。
如果名称已存在，
\lstinline{dirlink}
返回错误
\linerefs{kernel/fs.c:/Check.that.name.is.not.present/,/^..}/}。
主循环读取目录条目以查找未分配的条目。
当找到一个时，它提前停止循环
\linerefs{kernel/fs.c:/Look for an empty dirent/,/.*break/}，
将
\lstinline{off}
设置为可用条目的偏移量。
否则，循环以
\lstinline{off}
设置为
\lstinline{dp->size}
结束。
无论哪种方式，
\lstinline{dirlink}
然后通过写入偏移量
\lstinline{off}
将新条目添加到目录
\linerefs{kernel/fs.c:/if.writei/,/return -1/}。
%%
%%
%%
\section{代码：路径名}

路径名查找涉及对
\indexcode{dirlookup}
的连续调用，每个路径组件一个。
\lstinline{namei}
\lineref{kernel/fs.c:/^namei/}
评估
\lstinline{path}
并返回相应的
\lstinline{inode}。
函数
\indexcode{nameiparent}
是一个变体：它在最后一个元素之前停止，返回父目录的 inode 并将最终元素复制到
\lstinline{name}。
两者都调用通用函数
\indexcode{namex}
来完成实际工作。

\lstinline{namex}
\lineref{kernel/fs.c:/^namex/}
首先决定路径评估从哪里开始。
如果路径以斜杠开头，评估从根目录开始；
否则，从当前目录开始
\linerefs{kernel/fs.c:/..if.\*path.==....\)/,/idup/}。
然后它使用
\indexcode{skipelem}
依次考虑路径的每个元素
\lineref{kernel/fs.c:/while.*skipelem/}。
每次循环迭代必须在当前 inode
\lstinline{ip}
中查找
\lstinline{name}。
迭代从锁定
\lstinline{ip}
并检查它是否是目录开始。
如果不是，查找失败
\linerefs{kernel/fs.c:/^....ilock.ip/,/^....}/}。
（锁定
\lstinline{ip}
是必要的，不是因为
\lstinline{ip->type}
可能在脚下改变——它不能——而是因为直到
\indexcode{ilock}
运行，
\lstinline{ip->type}
不能保证已从磁盘加载。）
如果调用是
\indexcode{nameiparent}
且这是最后一个路径元素，循环提前停止，
根据
\lstinline{nameiparent}
的定义；
最终路径元素已复制到
\lstinline{name}，
因此
\indexcode{namex}
只需返回未锁定的
\lstinline{ip}
\linerefs{kernel/fs.c:/^....if.nameiparent/,/^....}/}。
最后，循环使用
\indexcode{dirlookup}
查找路径元素并通过设置
\lstinline{ip = next}
准备下一次迭代
\linerefs{kernel/fs.c:/^....if..next.*dirlookup/,/^....ip.=.next/}。
当循环用完路径元素时，它返回
\lstinline{ip}。

过程
\lstinline{namex}
可能需要很长时间才能完成：它可能涉及几个磁盘操作来读取路径中遍历的目录的 inode 和目录块（如果它们不在缓冲区缓存中）。Xv6 经过精心设计，以便如果一个内核对
\lstinline{namex}
的调用被磁盘 I/O 阻塞，另一个内核查找不同路径名可以并发进行。
\lstinline{namex}
单独锁定路径中的每个目录，以便不同目录中的查找可以并行进行。

这种并发引入了一些挑战。例如，当一个内核线程正在查找路径名时，另一个内核线程可能正在通过取消链接目录来更改目录树。一个潜在的风险是查找可能正在搜索已被另一个内核线程删除的目录，其块已被重用于另一个目录或文件。

Xv6 避免了这种竞争。例如，当在
\lstinline{namex}
中执行
\lstinline{dirlookup}
时，查找线程持有目录的锁，
\lstinline{dirlookup}
返回使用
\lstinline{iget}
获得的 inode。
\lstinline{iget}
增加 inode 的引用计数。只有在从
\lstinline{dirlookup}
接收到 inode 后，
\lstinline{namex}
才释放目录上的锁。现在另一个线程可以从目录中取消链接 inode，但 xv6 还不会删除 inode，因为 inode 的引用计数仍然大于零。

另一个风险是死锁。例如，
\lstinline{next}
在查找 "。" 时指向与
\lstinline{ip}
相同的 inode。
在释放
\lstinline{ip}
上的锁之前锁定
\lstinline{next}
将导致死锁。
为了避免这种死锁，
\lstinline{namex}
在获取
\lstinline{next}
上的锁之前解锁目录。
这里我们再次看到为什么
\lstinline{iget}
和
\lstinline{ilock}
之间的分离很重要。
%%
%%
%%
\section{文件描述符层}

Unix 接口的一个很酷的方面是 Unix 中的大多数资源都表示为文件，包括控制台、管道等设备，当然还有真正的文件。文件描述符层是实现这种统一性的层。

Xv6 给每个进程自己的打开文件表，或文件描述符，如我们在第~\ref{CH:UNIX} 章看到的。
每个打开的文件由
\indexcode{struct file}
\lineref{kernel/file.h:/^struct.file/}
表示，它是一个围绕 inode 或管道的包装器，加上一个 I/O 偏移量。
每次对
\indexcode{open}
的调用创建一个新的打开文件（一个新的
\lstinline{struct}
\lstinline{file}）：
如果多个进程独立打开同一个文件，不同的实例将具有不同的 I/O 偏移量。
另一方面，单个打开文件（相同的
\lstinline{struct}
\lstinline{file}）
可以在一个进程的文件表中出现多次，也可以在多个进程的文件表中出现。
如果一个进程使用
\lstinline{open}
打开文件，然后使用
\indexcode{dup}
创建别名或使用
\indexcode{fork}
与子进程共享，就会发生这种情况。
引用计数跟踪对特定打开文件的引用数量。
文件可以为读或写或两者打开。
\lstinline{readable}
和
\lstinline{writable}
字段跟踪这一点。

系统中的所有打开文件都保存在全局文件表
\indexcode{ftable}
中。
文件表具有分配文件
（\indexcode{filealloc}）、
创建重复引用
（\indexcode{filedup}）、
释放引用
（\indexcode{fileclose}）、
以及读写数据
（\indexcode{fileread}
和
\indexcode{filewrite}）的函数。

前三个遵循现在熟悉的形式。
\lstinline{filealloc}
\lineref{kernel/file.c:/^filealloc/}
扫描文件表以查找未引用的文件
（\lstinline{f->ref}
\lstinline{==}
\lstinline{0}）
并返回新引用；
\indexcode{filedup}
\lineref{kernel/file.c:/^filedup/}
递增引用计数；
\indexcode{fileclose}
\lineref{kernel/file.c:/^fileclose/}
递减它。
当文件的引用计数达到零时，
\lstinline{fileclose}
根据类型释放底层管道或 inode。

函数
\indexcode{filestat}、
\indexcode{fileread}、
和
\indexcode{filewrite}
实现
\indexcode{stat}、
\indexcode{read}、
和
\indexcode{write}
文件操作。
\lstinline{filestat}
\lineref{kernel/file.c:/^filestat/}
仅允许在 inode 上并调用
\indexcode{stati}。
\lstinline{fileread}
和
\lstinline{filewrite}
检查打开模式是否允许操作，然后
将调用传递给管道或 inode 实现。
如果文件表示 inode，
\lstinline{fileread}
和
\lstinline{filewrite}
使用 I/O 偏移量作为操作的偏移量，然后推进它
\linerefs{kernel/file.c:/readi/,/./}
\linerefs{kernel/file.c:/writei/,/./}。
管道没有偏移量的概念。
回想一下 inode 函数要求调用者处理锁定
\linerefs{kernel/file.c:/stati/-1,/iunlock/}
\linerefs{kernel/file.c:/readi/-1,/iunlock/}
\linerefs{kernel/file.c:/writei\(f/-1,/iunlock/}。
inode 锁定有一个方便的副作用，即读写偏移量原子地更新，以便多个同时写入同一文件不能覆盖彼此的数据，尽管它们的写入可能最终交错。
%%
%%
%%
\section{代码：系统调用}

有了低层提供的函数，大多数系统调用的实现是微不足道的
（参见
\fileref{kernel/sysfile.c}）。
有一些调用值得仔细研究。

函数
\indexcode{sys_link}
和
\indexcode{sys_unlink}
编辑目录，创建或删除对 inode 的引用。
它们是使用事务的强大功能的另一个好例子。
\lstinline{sys_link}
\lineref{kernel/sysfile.c:/^sys_link/}
首先获取其参数，两个字符串
\lstinline{old}
和
\lstinline{new}
\lineref{kernel/sysfile.c:/argstr.*old.*new/}。
假设
\lstinline{old}
存在且不是目录
\linerefs{kernel/sysfile.c:/namei.old/,/^..}/}，
\lstinline{sys_link}
递增其
\lstinline{ip->nlink}
计数。
然后
\lstinline{sys_link}
调用
\indexcode{nameiparent}
以查找
\lstinline{new}
的父目录和最终路径元素
\lineref{kernel/sysfile.c:/nameiparent.new/}
并创建一个新的目录条目指向
\lstinline{old}
的 inode
\lineref{kernel/sysfile.c:/\|\| dirlink/}。
新的父目录必须存在且与现有 inode 在同一设备上：
inode 号只在单个磁盘上具有唯一含义。
如果发生这样的错误，
\indexcode{sys_link}
必须返回并递减
\lstinline{ip->nlink}。

事务简化了实现，因为它需要更新多个磁盘块，但我们不必担心执行它们的顺序。它们要么全部成功，要么都不成功。
例如，没有事务，在创建链接之前更新
\lstinline{ip->nlink}
将使文件系统暂时处于不安全状态，中间的崩溃可能导致严重破坏。
有了事务，我们不必担心这个。

\lstinline{sys_link}
为现有 inode 创建新名称。
函数
\indexcode{create}
\lineref{kernel/sysfile.c:/^create/}
为新 inode 创建新名称。
它是三个文件创建系统调用的概括：
带
\indexcode{O_CREATE}
标志的
\indexcode{open}
创建新普通文件，
\indexcode{mkdir}
创建新目录，
\indexcode{mkdev}
创建设备文件。
像
\indexcode{sys_link}
一样，
\indexcode{create}
首先调用
\indexcode{nameiparent}
以获取父目录的 inode。
然后它调用
\indexcode{dirlookup}
检查名称是否已存在
\lineref{kernel/sysfile.c:/dirlookup.*\[^=\]=.0/}。
如果名称确实存在，
\lstinline{create}
的行为取决于它为哪个系统调用使用：
\lstinline{open}
与
\indexcode{mkdir}
和
\indexcode{mkdev}
具有不同的语义。
如果
\lstinline{create}
代表
\lstinline{open}
使用
（\lstinline{type}
\lstinline{==}
\indexcode{T_FILE}）
且存在的名称本身是普通文件，
则
\lstinline{open}
将其视为成功，
因此
\lstinline{create}
也视为成功
\lineref{kernel/sysfile.c:/^......return.ip/}。
否则，它是错误
\linerefs{kernel/sysfile.c:/^......return.ip/+1,/return.0/}。
如果名称尚不存在，
\lstinline{create}
现在使用
\indexcode{ialloc}
分配新 inode
\lineref{kernel/sysfile.c:/ialloc/}。
如果新 inode 是目录，
\lstinline{create}
用
\indexcode{.}
和
\indexcode{..}
条目初始化它。
最后，既然数据已正确初始化，
\indexcode{create}
可以将其链接到父目录
\lineref{kernel/sysfile.c:/if.dirlink/}。
\lstinline{create}，
像
\indexcode{sys_link}
一样，同时持有两个 inode 锁：
\lstinline{ip}
和
\lstinline{dp}。
没有死锁的可能性，因为 inode
\lstinline{ip}
是新分配的：系统中的其他进程不会持有
\lstinline{ip}
的锁然后试图锁定
\lstinline{dp}。

使用
\lstinline{create}，
很容易实现
\indexcode{sys_open}、
\indexcode{sys_mkdir}、
和
\indexcode{sys_mknod}。
\lstinline{sys_open}
\lineref{kernel/sysfile.c:/^sys_open/}
是最复杂的，因为创建新文件只是它能做的一小部分。
如果
\indexcode{open}
被传递
\indexcode{O_CREATE}
标志，它调用
\lstinline{create}
\lineref{kernel/sysfile.c:/create.*T_FILE/}。
否则，它调用
\indexcode{namei}
\lineref{kernel/sysfile.c:/if..ip.=.namei.path/}。
\lstinline{create}
返回锁定的 inode，但
\lstinline{namei}
不返回，因此
\indexcode{sys_open}
必须自己锁定 inode。
这提供了一个方便的地方来检查目录是否只打开用于读取而不是写入。
假设 inode 以某种方式获得，
\lstinline{sys_open}
分配文件和文件描述符
\lineref{kernel/sysfile.c:/filealloc.*fdalloc/}
然后填充文件
\linerefs{kernel/sysfile.c:/type.=.FD_INODE/,/writable/}。
请注意，没有其他进程可以访问部分初始化的文件，因为它只在当前进程的表中。

第~\ref{CH:SLEEP} 章在我们拥有文件系统之前检查了管道的实现。
函数
\indexcode{sys_pipe}
通过提供创建管道对的方式将该实现连接到文件系统。
它的参数是指向两个整数空间的指针，
它将在其中记录两个新的文件描述符。
然后它分配管道并安装文件描述符。
%%
%%  -------------------------------------------
%%
\section{现实世界}

真实操作系统中的缓冲区缓存比 xv6 的复杂得多，但它服务于相同的两个目的：缓存和同步对磁盘的访问。
Xv6 的缓冲区缓存，像 V6 的一样，使用简单的最近最少使用（LRU）逐出策略；
可以实现许多更复杂的策略，每种策略对某些工作负载好，对其他工作负载不太好。
更高效的 LRU 缓存将消除链表，
而是使用哈希表进行查找，使用堆进行 LRU 逐出。
现代缓冲区缓存通常与虚拟内存系统集成以支持内存映射文件。

Xv6 的日志系统效率低下。
提交不能与文件系统系统调用并发发生。
系统记录整个块，即使块中只有几个字节被更改。它执行同步日志写入，一次一个块，每个都可能需要整个磁盘旋转时间。真正的日志系统解决了所有这些问题。

日志不是提供崩溃恢复的唯一方法。早期文件系统在重启期间使用清除器（例如，Unix
\indexcode{fsck}
程序）来检查每个文件和目录以及块和 inode 空闲列表，寻找并解决不一致性。清除大型文件系统可能需要数小时，而且有些情况下无法以使原始系统调用原子化的方式解决不一致性。从日志中恢复要快得多，并使系统调用在面对崩溃时是原子的。

Xv6 使用与早期 Unix 相同的 inode 和目录的基本磁盘布局；
这种方案多年来一直非常持久。
BSD 的 UFS/FFS 和 Linux 的 ext2/ext3 使用本质上相同的数据结构。
文件系统布局中效率最低的部分是目录，它需要在每次查找时对所有磁盘块进行线性扫描。
当目录只有几个磁盘块时这是合理的，但对于包含许多文件的目录来说很昂贵。
Microsoft Windows 的 NTFS、macOS 的 HFS 和 Solaris 的 ZFS，仅举几例，将目录实现为磁盘上块的平衡树。
这很复杂，但保证了对数时间的目录查找。

Xv6 对磁盘故障很天真：如果磁盘操作失败，xv6 发生 panic。
这是否合理取决于硬件：
如果操作系统位于使用冗余来屏蔽磁盘故障的特殊硬件之上，也许操作系统很少看到故障，因此 panic 是可以的。
另一方面，使用普通磁盘的操作系统应该预期故障并更优雅地处理它们，
以便一个文件中一个块的丢失不会影响文件系统其余部分的使用。

Xv6 要求文件系统适合一个磁盘设备且大小不变。
随着大型数据库和多媒体文件推动存储需求越来越高，操作系统正在开发消除``每个文件系统一个磁盘''瓶颈的方法。
基本方法是将许多磁盘组合成一个逻辑磁盘。硬件解决方案如 RAID 仍然是最受欢迎的，但当前趋势是尽可能在软件中实现这种逻辑。
这些软件实现通常允许丰富的功能，如通过动态添加或删除磁盘来增长或缩小逻辑设备。
当然，可以动态增长或缩小的存储层需要可以执行相同操作的文件系统：xv6 使用的 inode 块的固定大小数组在此类环境中效果不佳。
将磁盘管理与文件系统分离可能是最简洁的设计，但两者之间的复杂接口导致一些系统，如 Sun 的 ZFS，将它们组合在一起。

Xv6 的文件系统缺乏现代文件系统的许多其他功能；例如，它缺乏对快照和增量备份的支持。

现代 Unix 系统允许使用与磁盘存储相同的系统调用来访问许多类型的资源：
命名管道、网络连接、远程访问的网络文件系统，以及监视和控制接口，如
\lstinline{/proc}。
与 xv6 的
\lstinline{if}
语句在
\indexcode{fileread}
和
\indexcode{filewrite}
中不同，这些系统通常给每个打开的文件一个函数指针表，
每个操作一个，
并调用函数指针来调用该 inode 的调用实现。
网络文件系统和用户级文件系统提供将这些调用转换为网络 RPC 的函数，并在返回之前等待响应。
%%
%%  -------------------------------------------
%%
\section{练习}

\begin{enumerate}

\item 为什么
\lstinline{balloc}
会发生 panic？
xv6 能恢复吗？

\item 为什么
\lstinline{ialloc}
会发生 panic？
xv6 能恢复吗？

\item 为什么
\lstinline{filealloc}
在文件用完时不发生 panic？
为什么这更常见因此值得处理？

\item 假设与
\lstinline{ip}
对应的文件在
\lstinline{sys_link}
对
\lstinline{iunlock(ip)}
和
\lstinline{dirlink}
的调用之间被另一个进程取消链接。
链接会正确创建吗？
为什么或为什么不？

\item
\lstinline{create}
进行四个函数调用（一个到
\lstinline{ialloc}
和三个到
\lstinline{dirlink}）
它要求成功。
如果任何一个失败，
\lstinline{create}
调用
\lstinline{panic}。
为什么这是可接受的？
为什么这四个调用都不能失败？

\item
\lstinline{sys_chdir}
在
\lstinline{iput(cp->cwd)}
之前调用
\lstinline{iunlock(ip)}，
后者可能试图锁定
\lstinline{cp->cwd}，
但将
\lstinline{iunlock(ip)}
推迟到
\lstinline{iput}
之后不会导致死锁。
为什么不呢？

\item 实现
\lstinline{lseek}
系统调用。支持
\lstinline{lseek}
还需要修改
\lstinline{filewrite}
以用零填充文件中的空洞，如果
\lstinline{lseek}
设置
\lstinline{off}
超过
\lstinline{f->ip->size.}

\item 将
\lstinline{O_TRUNC}
和
\lstinline{O_APPEND}
添加到
\lstinline{open}，
以便
\lstinline{>}
和
\lstinline{>>}
操作符在 shell 中工作。

\item 修改文件系统以支持符号链接。

\item 修改文件系统以支持命名管道。

\item 修改文件和 VM 系统以支持内存映射文件。

\end{enumerate}
